\chapter{Moving frame 与 Gauss--Codazzi--Ricci 方程}

前面所有曲率都属于流形自身的内在几何。若一个 Riemann 流形 $M^m$ 浸入更高维 Riemann 流形 $\overline M^{m+q}$，环境 Levi--Civita 联络沿 $M$ 的导数既有切向分量，也有法向分量。后者就是外在弯曲。最一般的兼容条件不是只适用于曲面的 Gauss--Codazzi 两式，而是任意余维的 Gauss、Codazzi 与 Ricci 三个曲率块。本章先在一般 $(m,q)$ 情形推导它们，再把 $M^2\subset\RR^3$ 作为余维一、平坦环境的特例。

\section{切向与法向联络的分块}

设
\[
F:M^m\longrightarrow\overline M^{m+q}
\]
为等距浸入。以后把 $TM$ 与其像识别。沿 $M$ 有正交分解
\[
T\overline M|_M=TM\oplus NM,
\]
其中 $NM$ 是法丛。环境度量与 Levi--Civita 联络分别记为 $\overline g$、$\overline\nabla$，诱导度量与 Levi--Civita 联络记为 $g$、$\nabla$。

对切向量场 $X,Y\in\Gamma(TM)$，把环境导数分解为
\begin{equation*}\boxed{
\overline\nabla_XY
=\nabla_XY+B(X,Y),}
\end{equation*}
其中
\[
B(X,Y):=(\overline\nabla_XY)^\perp\in\Gamma(NM)
\]
称为第二基本形式。环境联络无挠，且 $\nabla$ 也是无挠的，因此
\[
B(X,Y)-B(Y,X)
=(\overline\nabla_XY-\overline\nabla_YX-[X,Y])^\perp=0,
\]
故 $B$ 对称。

对法向场 $\xi\in\Gamma(NM)$，环境导数同样分解为
\begin{equation*}\boxed{
\overline\nabla_X\xi
=-A_\xi X+\nabla_X^\perp\xi.}
\end{equation*}
这里 $A_\xi:TM\to TM$ 是关于法向 $\xi$ 的 shape operator，$\nabla^\perp$ 是 normal connection。由 $X\overline g(Y,\xi)=0$ 与度量相容，
\[
0=\overline g(\overline\nabla_XY,\xi)
+\overline g(Y,\overline\nabla_X\xi),
\]
所以
\begin{equation}\label{eq:math-dg-shape-second-form-relation}
\boxed{
\overline g(B(X,Y),\xi)
=g(A_\xi X,Y).}
\end{equation}
由于 $B$ 对称，每个 $A_\xi$ 都关于 $g$ 自伴。

取适配正交标架
\[
\{e_i\}_{i=1}^m\subset TM,
\qquad
\{n_\alpha\}_{\alpha=1}^q\subset NM.
\]
写
\[
B(e_i,e_j)=h^\alpha{}_{ij}n_\alpha.
\]
则
\[
h^\alpha{}_{ij}=h^\alpha{}_{ji}.
\]
平均曲率向量定义为
\begin{equation*}\mathbf H
:=\frac1m\operatorname{tr}_gB.
\end{equation*}
它是法向量场，而不是内在标量曲率。

\section{Gauss--Codazzi--Ricci 的统一推导}

记 $\overline R$、$R$ 与 $R^\perp$ 分别为环境联络、切向 Levi--Civita 联络和 normal connection 的曲率，均采用
\[
R(X,Y)Z
=\nabla_X\nabla_YZ-\nabla_Y\nabla_XZ-\nabla_{[X,Y]}Z
\]
的符号约定。对第二基本形式定义协变导数
\begin{equation*}(\nabla_X^\perp B)(Y,Z)
:=\nabla_X^\perp(B(Y,Z))
-B(\nabla_XY,Z)-B(Y,\nabla_XZ).
\end{equation*}

\begin{derivation}[从环境曲率的切向/法向分块推出三条兼容方程]
先让环境曲率作用在切向量 $Z$。由 Gauss 公式，
\[
\overline\nabla_YZ
=\nabla_YZ+B(Y,Z).
\]
再沿 $X$ 求环境协变导数，并分别对切向与法向部分使用 Gauss--Weingarten 分解：
\[
\begin{aligned}
\overline\nabla_X\overline\nabla_YZ
={}&\nabla_X\nabla_YZ+B(X,\nabla_YZ)\\
&-A_{B(Y,Z)}X
+\nabla_X^\perp B(Y,Z).
\end{aligned}
\]
交换 $X,Y$，再减去 $\overline\nabla_{[X,Y]}Z$。切向部分得到
\[
(\overline R(X,Y)Z)^\top
=R(X,Y)Z
-A_{B(Y,Z)}X+A_{B(X,Z)}Y.
\]
与切向量 $W$ 做内积并用\eqnrefs{eq:math-dg-shape-second-form-relation}，
\[
\begin{aligned}
\overline g(\overline R(X,Y)Z,W)
={}&g(R(X,Y)Z,W)\\
&-\overline g(B(X,W),B(Y,Z))\\
&+\overline g(B(Y,W),B(X,Z)).
\end{aligned}
\]
移项后得到 Gauss 方程
\begin{equation}\label{eq:math-dg-gauss-general}
\boxed{
\begin{aligned}
g(R(X,Y)Z,W)
={}&\overline g(\overline R(X,Y)Z,W)\\
&+\overline g(B(X,W),B(Y,Z))\\
&-\overline g(B(X,Z),B(Y,W)).
\end{aligned}}
\end{equation}

同一个展开的法向部分为
\[
(\overline R(X,Y)Z)^\perp
=(\nabla_X^\perp B)(Y,Z)
-(\nabla_Y^\perp B)(X,Z),
\]
即 Codazzi 方程
\begin{equation}\label{eq:math-dg-codazzi-general}
\boxed{
(\overline R(X,Y)Z)^\perp
=(\nabla_X^\perp B)(Y,Z)
-(\nabla_Y^\perp B)(X,Z).}
\end{equation}

最后让环境曲率作用在法向场 $\xi$。由 Weingarten 公式，
\[
\overline\nabla_Y\xi
=-A_\xi Y+\nabla_Y^\perp\xi.
\]
再次沿 $X$ 求导并取法向部分：
\[
\begin{aligned}
(\overline\nabla_X\overline\nabla_Y\xi)^\perp
=-B(X,A_\xi Y)
+\nabla_X^\perp\nabla_Y^\perp\xi.
\end{aligned}
\]
交换 $X,Y$ 并减去 $[X,Y]$ 项，得到
\[
(\overline R(X,Y)\xi)^\perp
=R^\perp(X,Y)\xi
-B(X,A_\xi Y)+B(Y,A_\xi X).
\]
与另一法向场 $\eta$ 做内积，利用 shape operator 自伴性，
\[
\begin{aligned}
\overline g(\overline R(X,Y)\xi,\eta)
={}&\overline g(R^\perp(X,Y)\xi,\eta)\\
&-g(A_\eta X,A_\xi Y)
+g(A_\eta Y,A_\xi X)\\
={}&\overline g(R^\perp(X,Y)\xi,\eta)
-g([A_\xi,A_\eta]X,Y).
\end{aligned}
\]
因此 Ricci 方程为
\begin{equation}\label{eq:math-dg-ricci-general}
\boxed{
\overline g(R^\perp(X,Y)\xi,\eta)
=\overline g(\overline R(X,Y)\xi,\eta)
+g([A_\xi,A_\eta]X,Y).}
\end{equation}
交换子的次序已经由上述展开固定；若把 $[A_\xi,A_\eta]$ 反写，右端必须同时变号。

$\mathbb R^3$ 中的球面 $\mathbb S^2(R)$：法丛一维，$A_\xi=(1/R)\,\mathrm{id}$，$[A_\xi,A_\eta]=0$（一维法丛无交换子），Ricci 方程退化为 $\overline R$ 的切向分量。对 $\mathbb R^3$ 中环面 $T^2$：$A=0$（平面是极小且全测地），所有兼容方程平凡。对 $\mathbb R^4$ 中的 Clifford 环面 $T^2\subset\mathbb S^3\subset\mathbb R^4$：$A_{n_1}$ 和 $A_{n_2}$ 在切空间上分别为 $\pm J$（$J$ 为复结构），$[A_{n_1},A_{n_2}]=0$——法曲率消失，这是 Lawson 不等式的关键。
\end{derivation}

三条方程分别来自环境曲率的三个块：切--切块给 Gauss，切--法混合块给 Codazzi，法--法块给 Ricci。高余维时 shape operators 之间可能不交换，Ricci 方程正是 normal bundle 曲率与这种不交换性的联系。

\section{移动标架形式}

在适配正交标架 $\{e_i,n_\alpha\}$ 中，环境联络一形式分块为
\[
\overline\omega
=
\begin{pmatrix}
\omega^i{}_j & \overline\omega^i{}_\beta\\
\overline\omega^\alpha{}_j & \omega^\alpha{}_\beta
\end{pmatrix}.
\]
其中 $\omega^i{}_j$ 是 $TM$ 的 Levi--Civita 联络，$\omega^\alpha{}_\beta$ 是 normal connection，而混合块编码第二基本形式：
\begin{equation}\label{eq:math-dg-moving-frame-second-form}
\overline\omega^\alpha{}_i
=h^\alpha{}_{ij}\theta^j,
\qquad
\overline\omega^i{}_\alpha
=-h^\alpha{}_{ij}\theta^j.
\end{equation}
把这一分块代入环境第二结构方程
\[
\overline\Omega
=\dd\overline\omega+\overline\omega\wedge\overline\omega
\]
并分别读取切--切、切--法与法--法块，就得到\eqnrefs{eq:math-dg-gauss-general}、\eqnrefs{eq:math-dg-codazzi-general}与\eqnrefs{eq:math-dg-ricci-general}的外形式版本。移动标架法并不是另一套理论，而是同一分块曲率计算的紧凑表示。


上述三条兼容方程不仅是任何浸入都必须满足的必要条件；在平坦环境中，它们在适当拓扑假设下也是重构浸入的充分条件。这一事实称为子流形基本定理，它说明 Gauss--Codazzi--Ricci 并不是若干互不相关的恒等式，而是移动标架系统的完整可积条件。

设 $(M^m,g)$ 为单连通 Riemann 流形，$N\to M$ 是秩 $q$ 的 Euclidean 向量丛，带 metric connection $\nabla^\perp$，并给定对称双线性形式
\[
B\in\Gamma\bigl(\operatorname{Sym}^2T^*M\otimes N\bigr).
\]
若 $(g,\nabla^\perp,B)$ 满足 Euclidean 环境 $\overline R=0$ 时的 Gauss、Codazzi、Ricci 方程，则局部存在等距浸入
\[
F:M\longrightarrow\RR^{m+q}
\]
实现这些数据；在单连通情形可作全局重构，并且在固定初始点、初始适配标架以后唯一。去掉初始归一化后，唯一性只差一个 Euclidean 刚体运动。

\begin{derivation}[Gauss--Codazzi--Ricci 怎样成为浸入的可积条件]
在局部正交标架下，把 Levi--Civita 联络、normal connection 与第二基本形式组合成一个 $\mathfrak{so}(m+q)$-值一形式
\[
\mathcal A=
\begin{pmatrix}
\omega^i{}_j & -h^i{}_{\beta k}\theta^k\\
 h^\alpha{}_{jk}\theta^k & \omega^\alpha{}_{\beta}
\end{pmatrix}.
\]
这里上下混合块用环境度量识别，符号选择与\eqnrefs{eq:math-dg-moving-frame-second-form}一致。计算其曲率
\[
\mathcal F
:=\dd\mathcal A+\mathcal A\wedge\mathcal A.
\]
切--切块恰是 Gauss 方程的左减右，切--法块恰是 Codazzi 方程，法--法块恰是 Ricci 方程。因此
\begin{equation}
\text{Gauss+Codazzi+Ricci}
\quad\Longleftrightarrow\quad
\mathcal F=0.
\label{eq:math-dg-submanifold-flat-extended-connection}
\end{equation}
也就是说，全部兼容条件等价于这个扩展联络平坦。

在单连通区域上，平坦矩阵联络可积分为一个正交矩阵场 $Q:M\to O(m+q)$，使
\[
Q^{-1}\dd Q=\mathcal A
\]
（若采用相反的矩阵 convention，等价地写成 $\dd Q=Q\mathcal A$）。令 $Q_1,\ldots,Q_m$ 为前 $m$ 个列向量，并定义 $\RR^{m+q}$-值一形式
\[
\alpha:=\sum_{i=1}^m Q_i\,\theta^i.
\]
计算外微分：
\[
\dd\alpha
=\sum_i \dd Q_i\wedge\theta^i
+\sum_i Q_i\,\dd\theta^i.
\]
$\dd Q_i$ 的切向联络部分与无挠第一结构方程
\[
\dd\theta^i+\omega^i{}_j\wedge\theta^j=0
\]
恰好抵消；法向部分含
\[
h^\alpha{}_{ij}\theta^j\wedge\theta^i,
\]
又因 $h^\alpha{}_{ij}=h^\alpha{}_{ji}$ 对称而为零。故
\[
\dd\alpha=0.
\]
单连通性使闭的一形式 $\alpha$ 可全局积分，所以存在 $F:M\to\RR^{m+q}$ 满足
\[
\dd F=\alpha.
\]
因为 $Q_i$ 正交归一，
\[
\langle \dd F(X),\dd F(Y)\rangle
=\sum_i\theta^i(X)\theta^i(Y)=g(X,Y),
\]
故 $F$ 是等距浸入。再由 $\dd Q=Q\mathcal A$ 的混合块读出其第二基本形式与 normal connection 正是预先给定的 $B$ 与 $\nabla^\perp$。这完成了存在性重构的核心计算。

若两个重构具有相同初始点与初始适配标架，它们满足同一一阶系统，因此由唯一性完全一致；改变初始点和初始正交标架只对应环境中的平移与正交变换。于是唯一性恰好差一个 Euclidean 刚体运动。
\end{derivation}

子流形的 Gauss--Codazzi--Ricci 三方程由此等价于扩展联络的平坦性 $\mathcal F=\dd\mathcal A+\mathcal A\wedge\mathcal A=0$（式~\eqnrefs{eq:math-dg-submanifold-flat-extended-connection}）。

\section{Euclidean 三维空间中的曲面}

现在令 $m=2,q=1$，并取环境 $\overline M=\RR^3$ 的 Euclidean 度量。设 $(e_1,e_2,n)$ 为适配正交归一标架，余标架为 $(\theta^1,\theta^2,\theta^3)$。限制到曲面后
\[
\theta^3=0.
\]
环境第一结构方程的法向分量给出
\begin{equation*}\omega^3{}_1\wedge\theta^1
+\omega^3{}_2\wedge\theta^2=0.
\end{equation*}
写
\[
\omega^3{}_1=h_{11}\theta^1+h_{12}\theta^2,
\qquad
\omega^3{}_2=h_{21}\theta^1+h_{22}\theta^2,
\]
代回上式得到
\begin{equation*}h_{12}=h_{21},
\end{equation*}
这正是一般结论“$B$ 对称”在二维余维一移动标架中的分量形式。第二基本形式写成
\begin{equation*}II=h_{ij}\theta^i\theta^j.
\end{equation*}
按本书约定，shape operator $A$ 满足
\[
\overline\nabla_Xn=-AX,
\qquad
II(X,Y)=g(AX,Y).
\]
Gauss--Weingarten 公式于是成为
\begin{equation*}\boxed{
\begin{aligned}
\overline\nabla_XY&=\nabla_XY+II(X,Y)n,\\
\overline\nabla_Xn&=-AX.
\end{aligned}}
\end{equation*}

因为 Euclidean 环境平坦，$\overline R=0$。Gauss 方程在正交基 $e_1,e_2$ 上给出
\[
K
=g(R(e_1,e_2)e_2,e_1)
=h_{11}h_{22}-h_{12}^2,
\]
即
\begin{equation*}\boxed{K=\det A.}
\end{equation*}
若 $\kappa_1,\kappa_2$ 是主曲率，则
\[
K=\kappa_1\kappa_2,
\qquad
\mathsf H=\frac12(\kappa_1+\kappa_2),
\]
其中 $\mathsf H$ 表示有向标量平均曲率，平均曲率向量为 $\mathbf H=\mathsf H n$。采用不同法向会使 $\kappa_i$ 与 $\mathsf H$ 同时变号，但 $K$ 不变。

Codazzi 方程在平坦环境中化为
\[
(\nabla_X^\perp B)(Y,Z)
=(\nabla_Y^\perp B)(X,Z).
\]
余维一时 $\nabla^\perp n=0$，若定义
\[
Dh_{ij}
:=\dd h_{ij}-h_{kj}\omega^k{}_i-h_{ik}\omega^k{}_j
=h_{ij;k}\theta^k,
\]
则移动标架的两个独立方程可写为
\begin{equation*}\dd\omega^3{}_1+\omega^3{}_2\wedge\omega^2{}_1=0,
\end{equation*}
\begin{equation*}\dd\omega^3{}_2+\omega^3{}_1\wedge\omega^1{}_2=0,
\end{equation*}
并等价于分量式
\begin{equation*}\boxed{h_{ij;k}=h_{ik;j}.}
\end{equation*}
余维一只有一个 shape operator，所以 Ricci 方程中的交换子自动为零；normal bundle 也是一维，Ricci 方程不再给出额外独立条件。


对定向 Euclidean hypersurface
\[
F:M^m\to\RR^{m+1},
\]
单位法向本身定义 Gauss map
\begin{equation*}\mathcal N:M\longrightarrow S^m,
\qquad
\mathcal N(p)=n(p).\end{equation*}
由 Weingarten 公式，任意 $X\in T_pM$ 都有
\begin{equation}
\boxed{
\dd\mathcal N_p(X)
=\overline\nabla_Xn
=-A_pX.
}
\label{eq:math-dg-gauss-map-differential}
\end{equation}
所以 shape operator 不是抽象定义出来的矩阵：它就是 Gauss map 的负微分。其本征方向是法向变化最快的正交方向，主曲率则是这些方向上的有符号伸缩率。

把单位球面的标准度量记为 $g_{S^m}$，则 Gauss map 拉回度量为
\begin{equation*}(\mathcal N^*g_{S^m})(X,Y)
=g(AX,AY).\end{equation*}
右边常称为第三基本形式。另一方面，Gauss map 的有向 Jacobian 为
\begin{equation*}\det(\dd\mathcal N)
=(-1)^m\det A.\end{equation*}
对曲面 $M^2\subset\RR^3$，$\det A=K$，故 Gaussian curvature 恰是 Gauss map 的局部面积伸缩率。$K$ 的“内在性”和法向映射的“外在定义”在这里并不矛盾：Gauss 方程保证二者数值一致。

shape operator 还精确控制平行超曲面的首次奇性。定义
\[
F_s(p):=F(p)+s\,n(p).
\]
对 $X\in T_pM$，由\eqnrefs{eq:math-dg-gauss-map-differential}
\begin{equation*}\dd F_s(X)
=(I-sA)X.\end{equation*}
因此只要 $I-sA$ 可逆，$F_s$ 仍是局部浸入；在主方向 $e_i$ 上，
\[
\dd F_s(e_i)=(1-s\kappa_i)e_i.
\]
当 $s=1/\kappa_i$ 时该方向退化，这就是 focal point 出现的外在版本。

在未退化区域内，法向仍可选为同一个 $n$。由
\[
-A_s\bigl((I-sA)X\bigr)
=\overline\nabla_Xn
=-AX
\]
得到
\begin{equation*}\boxed{
A_s=A(I-sA)^{-1},
\qquad
\kappa_i(s)=\frac{\kappa_i}{1-s\kappa_i}.
}\end{equation*}
同时面积元按
\begin{equation}
\operatorname{vol}_{g_s}
=|\det(I-sA)|\,\operatorname{vol}_g
=\prod_{i=1}^m|1-s\kappa_i|\,\operatorname{vol}_g
\label{eq:math-dg-parallel-area-jacobian}
\end{equation}
变化。展开这个乘积就会依次出现平均曲率、第二基本形式的初等对称多项式以及 Gauss--Kronecker curvature。于是 tube formula 与曲率积分的系数，本质上都来自同一个 determinant。

\begin{derivation}[平行曲面的面积一阶变化如何恢复平均曲率]
由\eqnrefs{eq:math-dg-parallel-area-jacobian}，平行曲面的面积元为
\begin{equation*}
\operatorname{vol}_{g_s}=|\det(I-sA)|\,\operatorname{vol}_g
=\prod_{i=1}^m|1-s\kappa_i|\,\operatorname{vol}_g,
\end{equation*}
其中 $\kappa_1,\ldots,\kappa_m$ 是 shape operator $A$ 的本征值即主曲率。在 $s=0$ 附近假设尚未遇到 focal point，即 $1-s\kappa_i>0$ 对所有 $i$ 成立，可去掉绝对值，
\begin{equation*}
\operatorname{vol}_{g_s}=\det(I-sA)\,\operatorname{vol}_g
=\prod_{i=1}^m(1-s\kappa_i)\,\operatorname{vol}_g.
\end{equation*}
把行列式按特征多项式在 $s=0$ 处 Taylor 展开。由 $\det(I-sA)$ 是 $s$ 的 $m$ 次多项式，常数项为 $\det I=1$，一次项为 $-s\operatorname{tr}A$（特征多项式系数与初等对称多项式的关系，一次项系数恰为负迹），故
\begin{equation*}
\det(I-sA)=1-s\operatorname{tr}A+O(s^2).
\end{equation*}
等价地，对乘积形式逐项展开，
\begin{equation*}
\prod_{i=1}^m(1-s\kappa_i)
=1-s\sum_{i=1}^m\kappa_i+O(s^2)
=1-s\operatorname{tr}A+O(s^2),
\end{equation*}
其中用了 $\operatorname{tr}A=\sum_i\kappa_i$。由平均曲率定义 $\mathsf H=\frac1m\operatorname{tr}A=\frac1m\sum_i\kappa_i$，
\begin{equation*}
\operatorname{tr}A=m\mathsf H
\end{equation*}
（余维一时 $\mathbf H=\mathsf H n$）。代回面积元展开，
\begin{equation*}
\operatorname{vol}_{g_s}=(1-sm\mathsf H+O(s^2))\,\operatorname{vol}_g.
\end{equation*}
对 $s$ 求导并在 $s=0$ 取值，
\begin{align*}
\left.\frac{\dd}{\dd s}\operatorname{vol}_{g_s}\right|_{s=0}
&=\left.\frac{\dd}{\dd s}\bigl[(1-sm\mathsf H+O(s^2))\,\operatorname{vol}_g\bigr]\right|_{s=0}\\
&=\bigl[-m\mathsf H+O(s)\bigr]\Big|_{s=0}\operatorname{vol}_g\\
&=-m\mathsf H\,\operatorname{vol}_g.
\end{align*}
这正是后面\eqnrefs{eq:math-dg-first-variation-submanifold-volume}在常法向速度 $V=n$ 情形的局部版本。换言之，平均曲率既是 shape operator 的迹，也是平行法向推进时面积元的一阶响应。

高阶项与初等对称多项式。继续展开到二阶，
\begin{equation*}
\det(I-sA)=1-s\sigma_1+\frac{s^2}{2}(\sigma_1^2-\sigma_2)+O(s^3),
\end{equation*}
其中 $\sigma_k=\sum_{i_1<\cdots<i_k}\kappa_{i_1}\cdots\kappa_{i_k}$ 是主曲率的第 $k$ 个初等对称多项式。$\sigma_2$ 与 Gauss 曲率 $K=\det A=\sigma_m$ 通过 Newton 恒等式相关。二阶项
\begin{equation*}
\frac12\left.\frac{\dd^2}{\dd s^2}\operatorname{vol}_{g_s}\right|_{s=0}
=\frac12(\sigma_1^2-\sigma_2)\operatorname{vol}_g
=\frac{m^2\mathsf H^2-|A|^2}{2}\operatorname{vol}_g
\end{equation*}
给出二阶变分，是研究 minimal 子流形稳定性的关键。当 $m^2\mathsf H^2=|A|^2$ 时二阶项消失，对应全测地子流形。

管状邻域与 Weyl 公式。对 $\RR^n$ 中紧子流形 $M^m$，管状邻域体积为
\begin{equation*}
\operatorname{Vol}(T_r(M))
=\omega_{n-m}\int_M\int_0^r\prod_{i=1}^m(1-s\kappa_i)\,\dd s\,\operatorname{vol}_g,
\end{equation*}
其中 $\omega_{n-m}$ 是单位 $(n-m)$-球体积。展开后被积函数是 $\sigma_k$ 的组合，对偶数余维给出 Weyl 管公式
\begin{equation*}
\operatorname{Vol}(T_r(M))
=\sum_{k=0}^{\lfloor(n-m)/2\rfloor}\frac{\omega_{n-m-2k}r^{n-m-2k}}{(n-m-2k)\cdots(n-m)}\int_M\sigma_{2k}\,\operatorname{vol}_g.
\end{equation*}
只有偶数阶 $\sigma_{2k}$ 出现，反映管体积对法向反射对称。这些积分是 Lipschitz--Killing 曲率，与热核渐近展开系数一致。

举例：球面与圆柱。对 $S^m(r)\subset\RR^{m+1}$，主曲率全为 $\kappa_i=1/r$，平均曲率 $\mathsf H=1/r$。平行曲面 $S^m(r+s)$ 面积元为 $(1+s/r)^m\operatorname{vol}_g$，一阶项 $m/r\cdot s=m\mathsf H s$ 与公式一致。对圆柱 $S^1(r)\times\RR\subset\RR^3$，主曲率为 $(1/r,0)$，$\mathsf H=1/(2r)$，平行曲面面积元为 $(1+s/r)\operatorname{vol}_g$，一阶项 $s/r=2\mathsf H s$，验证余维一情形。

膜与肥皂膜。极小曲面方程 $\mathsf H=0$ 是肥皂膜的平衡条件，对应面积泛函的临界点。平行曲面公式给出膜在法向小扰动下面积的一阶响应：$\mathsf H=0$ 时一阶项消失，膜处于平衡；二阶项 $\frac12\int_M(m^2\mathsf H^2-|A|^2)\operatorname{vol}_g=-\frac12\int_M|A|^2\operatorname{vol}_g$ 在非全测地时为负，反映膜在大多数情形下不稳定，需要边界约束维持。这与 Plateau 问题的解存在性密切相关。
\end{derivation}

外在几何真正进入变分问题时，第二基本形式不再只是兼容条件中的系数，而直接控制面积的一阶与二阶变化。设
\[
F_t:M^m\to\overline M^{m+q}
\]
是一族浸入，变分速度为
\[
V:=\partial_tF_t|_{t=0}.
\]
把 $V=V^\top+V^\perp$ 分解成切向与法向部分。切向部分只对应参数重标记；内在面积真正改变的主项来自 $V^\perp$。对紧支撑变分，一阶体积公式为
\begin{equation}
\boxed{
\left.\frac{\dd}{\dd t}\operatorname{Vol}(F_t(M))\right|_{t=0}
=-m\int_M
\overline g(\mathbf H,V^\perp)\,\operatorname{vol}_g
}.
\label{eq:math-dg-first-variation-submanifold-volume}
\end{equation}
若有边界，还要加上由 $V^\top$ 产生的边界通量项。因此
\[
\boxed{\mathbf H=0}
\]
恰是对所有紧支撑法向变分面积临界的 Euler--Lagrange 方程；这样的浸入称为 minimal submanifold。

\begin{derivation}[第二基本形式怎样进入面积一阶变分]
取局部坐标 $x^i$，诱导度量为
\[
g_{ij}(t)
=\overline g(\partial_iF_t,\partial_jF_t).
\]
时间求导并用环境联络度量相容性：
\[
\dot g_{ij}
=\overline g(\overline\nabla_iV,e_j)
+\overline g(e_i,\overline\nabla_jV).
\]
因此
\[
\frac12\operatorname{tr}_g\dot g
=g^{ij}\overline g(\overline\nabla_iV,e_j).
\]
将 $V=V^\top+V^\perp$ 分解。切向部分给出
\[
g^{ij}g(\nabla_iV^\top,e_j)
=\operatorname{div}_M V^\top.
\]
对法向部分，利用
\[
0=e_i\,\overline g(V^\perp,e_j)
=\overline g(\overline\nabla_iV^\perp,e_j)
+\overline g(V^\perp,\overline\nabla_ie_j)
\]
得到
\[
\overline g(\overline\nabla_iV^\perp,e_j)
=-\overline g(V^\perp,B_{ij}).
\]
收缩后
\[
\frac12\operatorname{tr}_g\dot g
=\operatorname{div}_M V^\top
-m\overline g(\mathbf H,V^\perp).
\]
再用体积形式的一般变分公式
\[
\partial_t\operatorname{vol}_g
=\frac12\operatorname{tr}_g\dot g\,\operatorname{vol}_g
\]
并对散度项积分，即得到\eqnrefs{eq:math-dg-first-variation-submanifold-volume}。

意义。一阶变分公式 $\delta\operatorname{Area}=-\int_M m\,\overline g(\mathbf H,V^\perp)\dd\operatorname{vol}_g$ 说明面积泛函的 Euler--Lagrange 方程是平均曲率 $\mathbf H=0$——这就是极小子流形方程。切向变分只产生散度项（边界项），不改变面积；法向变分通过平均曲率与面积变化耦合。物理上，肥皂膜在表面张力下平衡时正是极小曲面：$\mathbf H=0$ 是 Young--Laplace 压力平衡在零压差时的特例。

 (1) $\mathbb R^3$ 中的平面：第二基本形式 $B_{ij}=0$，故 $\mathbf H=0$，面积在任意紧支撑变分下稳定。(2) $\mathbb R^3$ 中的球面 $\mathbb S^2(R)$：$B_{ij}=(1/R)g_{ij}$，$\mathbf H=(1/R)n\neq0$，球面不是极小曲面——缩小半径可减少面积。(3) 悬链面 $z=\cosh^{-1}(r)$：可验证 $\mathbf H=0$，这是 $\mathbb R^3$ 中唯一的非平凡旋转极小曲面，对应于两圆环间肥皂膜的平衡形状。
\end{derivation}

对余维一 minimal hypersurface，取单位法向 $n$ 与法向变分
\[
V=\phi n.
\]
在 $\mathbf H=0$ 的临界点，一阶变分消失，二阶变分由 stability quadratic form 控制：
\begin{equation}
\boxed{
\delta^2\operatorname{Vol}[\phi]
=
\int_M
\left(
|\nabla\phi|^2
-\bigl(|A|^2+\overline{\operatorname{Ric}}(n,n)\bigr)\phi^2
\right)\operatorname{vol}_g
}.
\label{eq:math-dg-minimal-second-variation}
\end{equation}
这与前面测地线的 index form 具有完全相同的结构：梯度项提供正的``弹性''，曲率势项可能制造负方向。事实上，一维 minimal submanifold 就是测地线，而\eqnrefs{eq:math-dg-minimal-second-variation}正退化为测地线 index form。

若\eqnrefs{eq:math-dg-minimal-second-variation}对所有紧支撑 $\phi$ 都非负，称 minimal hypersurface stable。因而稳定性问题被化成一个 Schrödinger 型二次型的谱下界问题；前面 Hilbert 空间、自伴算子与 Friedrichs 实现的理论在这里重新出现。几何变分、谱理论和 PDE 并不是三块独立知识，而是在这一 Jacobi quadratic form 上闭合。


把分部积分显式做出来，二阶变分可写成
\begin{equation*}\delta^2\operatorname{Vol}[\phi]
=\int_M\phi\,L_J\phi\,\operatorname{vol}_g,
\qquad
L_J:=-\Delta_{\mathrm{LB}}
-\left(|A|^2+\overline{\operatorname{Ric}}(n,n)\right).\end{equation*}
这里 $L_J$ 是与当前 Laplace--Beltrami 符号约定相容的 self-adjoint Jacobi operator。稳定性等价于 $L_J$ 的 quadratic form 非负；Morse index 则等于其负本征值总数（计重数）。这与测地线的 Morse index theorem 完全平行，只是普通 Sturm--Liouville 算子被 hypersurface 上的 Schrödinger operator 取代。

Euclidean minimal hypersurface 还有一个重要的二阶恒等式，把第二基本形式本身变成椭圆方程的未知量。若 $\overline M=\RR^{m+1}$ 且 $\mathsf H=0$，则 Codazzi 与 Gauss 方程推出 Simons identity
\begin{equation}
\boxed{
\Delta_{\mathrm{LB}}h_{ij}
+|A|^2h_{ij}=0.
}
\label{eq:math-dg-simons-hessian}
\end{equation}
与 $h_{ij}$ 收缩，并使用
\[
\frac12\Delta_{\mathrm{LB}}|A|^2
=\langle\Delta_{\mathrm{LB}}A,A\rangle
+|\nabla A|^2,
\]
得到标量形式
\begin{equation*}\boxed{
\frac12\Delta_{\mathrm{LB}}|A|^2
=|\nabla A|^2-|A|^4.
}\end{equation*}
它说明 minimal 条件并没有让外在曲率变成 harmonic；$|A|^4$ 是一个临界的非线性反应项，恰好与梯度耗散竞争。

\begin{derivation}[Simons identity 的结构来源]
在 Euclidean hypersurface 上，Codazzi 给出
\[
\nabla_k h_{ij}=\nabla_i h_{kj}.
\]
因此
\[
\Delta_{\mathrm{LB}}h_{ij}
=\nabla^k\nabla_kh_{ij}
=\nabla^k\nabla_i h_{kj}.
\]
交换两次协变导数时产生 Riemann curvature 对两个切指标的作用；再用 Gauss 方程把所有内在 Riemann tensor 写成 $h*h$。另一方面，minimal 条件
\[
\operatorname{tr}A=0
\]
使 $\nabla_i\nabla_j(\operatorname{tr}A)$ 项完全消失。剩余的所有三次项整理后恰为
\[
-|A|^2h_{ij},
\]
从而得到\eqnrefs{eq:math-dg-simons-hessian}。最后对
\[
|A|^2=h_{ij}h^{ij}
\]
应用乘积法则：
\[
\frac12\Delta_{\mathrm{LB}}|A|^2
=|\nabla A|^2
+\langle\Delta_{\mathrm{LB}}A,A\rangle
=|\nabla A|^2-|A|^4.
\]
所以 Simons identity 的三项来源分别非常清楚：Codazzi 负责交换一个导数，Gauss 负责把交换子曲率变成 $A^2$，minimality 负责杀掉 mean-curvature Hessian。

交换子计算的细节。由 Riemann 曲率作用，
\begin{equation*}
(\nabla^k\nabla_i-\nabla_i\nabla^k)h_{kj}
=R^k{}_\ell h_{\ell j}+R_j{}_\ell h_{k\ell}
-R^\ell{}_{kj}h_{i\ell},
\end{equation*}
其中 $R_{ij}$ 是 Ricci 张量。代入 Gauss 方程 $R_{ij}=h_{i\ell}h_j{}^\ell-Hh_{ij}$（$H=\operatorname{tr}A$ 为平均曲率），并使用 minimality $H=0$，得
\begin{equation*}
R^k{}_\ell h_{\ell j}=h^k{}_\ell h^\ell{}_m h^m{}_j=|A|^2h^k{}_j,
\end{equation*}
最后一步用到 $h$ 是对称的且 $h^k{}_\ell h^\ell{}_m=|A|^2\delta^k_m$ 在主曲率坐标系下成立。其他两项类似计算，合并后给出 $-|A|^2h_{ij}$。

Bernstein 定理的 Simons 证明。对 minimal graph $u:\Omega\subset\RR^n\to\RR$，Simons identity 结合 Bochner 技巧给出
\begin{equation*}
\frac12\Delta_{\mathrm{LB}}|A|^2\geq-|A|^4+\frac{2}{n}|A|^4
=\left(\frac{2}{n}-1\right)|A|^4.
\end{equation*}
当 $n\leq7$ 时 $\frac{2}{n}-1<0$，配合 $|A|^2$ 的有界性给出 $|A|^2\equiv0$，即 $u$ 是线性函数。这就是 Bernstein 定理在 $n\leq7$ 的 Simons 证明：整 minimal graph 必须是超平面。$n=8$ 时 Bombieri--De Giorgi--Giusti 给出反例，临界维数正是 Simons 不等式系数变号的边界。

Chern--Simons 形式的关联。对三维流形上的 $SU(2)$ 联络，Chern--Simons 形式
\begin{equation*}
\operatorname{CS}(A)=\operatorname{tr}\left(A\wedge\dd A+\frac23A\wedge A\wedge A\right)
\end{equation*}
的变分给出 $F\wedge F$，与 Simons identity 中 $|A|^4$ 项有类似的"曲率平方"结构。这是 Chern--Simons 拓扑场论与极小曲面理论在变分结构上的深层联系，也是 Witten 对 Jones 多项式几何解释的出发点。

举例：Clifford 极小环面。在 $S^3\subset\RR^4$ 中，Clifford 极小环面
\begin{equation*}
T=\left\{\frac{1}{\sqrt2}(\cos u,\sin u,\cos v,\sin v):u,v\in[0,2\pi]\right\}
\end{equation*}
是 $S^3$ 中唯一非全测地紧致极小曲面（Chern do Carmo--Lawson 给出）。其 $|A|^2=2$ 处处为常数，Simons identity 给出 $\Delta_{\mathrm{LB}}|A|^2=0$，对应 $|\nabla A|^2=|A|^4=4$，验证 $A$ 不能平行：极小嵌入与全测地之间存在精确的非线性平衡。

弦论中的极小曲面。弦论中开弦的世界面是 Dirichlet 边界条件下的极小曲面，闭弦的世界面是体积泛函的临界点。Simons identity 给出世界面方程在量子涨落下稳定性的判据：$|A|^2$ 满足的椭圆不等式控制弦世界面在背景时空中的涨落尺度，与弦论中 $\alpha'$ 修正的几何展开直接相关。
\end{derivation}


对图形浸入，minimal 方程可以直接写成一个拟线性椭圆 PDE。令
\[
F(x)=(x,u(x)),
\qquad x\in\Omega\subset\RR^m.
\]
诱导度量满足
\[
g_{ij}=\delta_{ij}+u_i u_j,
\qquad
\det g=1+|\nabla u|^2.
\]
因此面积泛函为
\begin{equation*}\mathcal A[u]
=\int_\Omega\sqrt{1+|\nabla u|^2}\,\dd^m x.\end{equation*}
做紧支撑变分 $u\mapsto u+\varepsilon\eta$，有
\[
\begin{aligned}
\delta\mathcal A[u;\eta]
&=\int_\Omega
\frac{\nabla u\cdot\nabla\eta}
{\sqrt{1+|\nabla u|^2}}\,\dd^mx\\
&=-\int_\Omega
\eta\,
\operatorname{div}\left(
\frac{\nabla u}{\sqrt{1+|\nabla u|^2}}
\right)\dd^mx.
\end{aligned}
\]
所以 minimal graph equation 是
\begin{equation*}\boxed{
\operatorname{div}\left(
\frac{\nabla u}{\sqrt{1+|\nabla u|^2}}
\right)=0
}.\end{equation*}
在小斜率 $|\nabla u|\ll1$ 下展开
\[
\frac{1}{\sqrt{1+|\nabla u|^2}}
=1-\frac12|\nabla u|^2+O(|\nabla u|^4),
\]
最低阶退化为
\[
\Delta u=0.
\]
因此 harmonic graph 是 minimal graph 的线性化，而不是与其无关的偶然近似。

面积最速下降则给出 mean curvature flow。若选择纯法向速度
\[
V=m\mathbf H,
\]
则由一阶变分公式
\[
\frac{\dd}{\dd t}\operatorname{Vol}(M_t)
=-m^2\int_{M_t}|\mathbf H|^2\,\operatorname{vol}_{g_t}\le0.
\]
这说明 mean curvature flow 与 Ricci flow 在逻辑上不同：前者演化的是浸入并以外在平均曲率降低面积，后者直接演化内在度量。两者都可看作几何泛函的梯度型演化，但状态空间和 gauge 结构不同。

\begin{exbox}[单位球面与圆柱面]
单位球面取外法向 $n=r$。对切向量 $X$，
\[
\overline\nabla_Xn=X.
\]
由约定 $\overline\nabla_Xn=-AX$，有 $A=-I$，故
\[
\kappa_1=\kappa_2=-1,
\qquad
K=1,
\qquad
\mathsf H=-1.
\]

半径 $R$ 的圆柱面取外法向时，圆周方向主曲率为 $-1/R$，轴向主曲率为零，因此
\[
K=0,
\qquad
\mathsf H=-\frac{1}{2R}.
\]
圆柱具有非零外在弯曲却有零 Gaussian curvature，直接说明平均曲率依赖嵌入，而 $K$ 是内在度量可以检测的曲率。
\end{exbox}

一般余维的三条兼容方程到此才在 $q=1$、$\overline R=0$ 的条件下退化为熟悉的曲面公式。这样的顺序保证 Gauss--Codazzi 不是从三维曲面的特殊计算“推广”出来，而是环境联络曲率分块的一般结论。
